% Cote 119: offprints, reports, courses, CV and letters (typed and printed part): archive cote 119, site page 64 (PDF page 65).
% offprint. Draft compared: none.
% AI-transcribed by Claude Opus 5.5 (claude-opus-5-5, adaptive thinking, effort high), Message Batches API, batch r119_full (msgbatch_01DwZNm5Y7YsKBWyAN4sFs8w), 2026-11-13; not yet checked by eye.
\scanpage{119/64}
\mspage{( 4 )}

Théorème 3.3. --- Soit $C$ une catégorie tannakienne de lien L. Alors,
si $C$ est polarisable [i. e. Pol $(C) \neq \emptyset$], toute catégorie tannakienne de lien L
est polarisable. De plus, l'application
\[ (3.3.1) \qquad \mathrm{H}^2(\mathbf{R}, \mathrm{L}) \rightarrow \mathrm{H}^2(\mathbf{R}, \mathrm{Z}), \qquad C \mapsto \varepsilon_C. \]
qu'on en déduit [voir ($^1$), 1.6] est compatible avec les actions de $\mathrm{H}^2(\mathbf{R}, \mathrm{Z})$
[voir ($^5$), 3.3, où on munit $\mathrm{H}^2(\mathbf{R}, \mathrm{L})$ d'une structure de pseudo-torseur
sous $\mathrm{H}^2(\mathbf{R}, \mathrm{Z})$], et en particulier est une bijection.

Corollaire 3.4 (Théorème de rigidité). --- Si $C$, $C'$ sont des catégories
tannakiennes de lien L, munies de polarisations $\pi$, $\pi'$ avec $\varepsilon(\pi) = \varepsilon(\pi')$,
il existe une équivalence de L-catégories tannakiennes $C \simeq C'$ unique à isomor-%
phisme (non unique) près respectant les polarisations données.

Définition 3.5. --- Un lien L sur $\mathbf{R}$ (défini localement par un groupe)
algébrique affine est dit polarisable s'il existe une catégorie tannakienne
liée par L qui soit polarisable.

4. Détermination des liens polarisables.

Théorème 4.1. --- Soit L un lien algébrique affine qui est, soit connexe,
soit abélien. Alors L est polarisable si et seulement si L est le lien d'un
$\mathbf{R}$-groupe compact. De plus, dans ce cas, l'unique classe $e \in \mathrm{H}^2(\mathbf{R}, \mathrm{L})$ envoyée
par l'application (3.3.1) en 0 est neutre.

Corollaire 4.2. --- Soit $C$ une catégorie tannakienne dont le lien L est
comme ci-dessus. Alors $C$ possède une polarisation symétrique si et seulement
s'il existe un foncteur fibre $\omega : C \rightarrow Mod\,f(\mathbf{R})$ avec $Aut^{\otimes}(\omega)$ compact.
De plus, la correspondance qui, à un tel $\omega$, associe la polarisation par les
formes $\varphi : \mathrm{V} \otimes \mathrm{V} \rightarrow 1$ telles que $\omega(\varphi)$ soit symétrique définie positive induit
une bijection de l'ensemble des classes d'isomorphisme de tels foncteurs fibre
sur l'ensemble $\mathrm{Pol}_1(C)$ des polarisations symétriques sur $C$.

Corollaire 4.3. --- Soit G un $\mathbf{R}$-groupe algébrique affine connexe ou
abélien. Alors, $Rep_0(\mathrm{G})$ est polarisable si et seulement si G est hodgien (2.5)
et, dans ce cas, toutes les polarisations sont hodgiennes.

Question 4.4. --- A-t-on des énoncés analogues sans supposer L connexe
ou abélien ? Le premier cas à décider est celui des $\mathbf{R}$-groupes finis (non
constants).

(*) Séance du 29 novembre 1971.
($^1$) N. Saavedra, Comptes rendus, 272, série A, 1971, p. 398.
($^2$) Pour des généralités concernant les catégories tannakiennes, on pourra se reporter
à un ouvrage en cours de préparation, à paraître dans les Lecture Notes in Mathematics,
Springer-Verlag.
($^3$) P. Deligne, Travaux de Griffiths, Séminaire Bourbaki, 376 (mai-juin 1970).
($^4$) P. Deligne, La conjecture de Weil pour les surfaces $\mathrm{K}_3$ (à paraître dans Inventiones
Mathematicae).
($^5$) J. Giraud, Cohomologie non abélienne, Springer-Verlag, Berlin, 1971.

Institut des Hautes Études scientifiques,
35, route de Chartres,
91-Bures-sur-Yvette, Essonne.

184937. --- Imp. Gauthier-Villars. --- 55, Quai des Grands-Augustins, Paris (6$^{\mathrm{e}}$)
Imprimé en France

\mspage{\hand{64}}
